\documentclass[a4paper,11pt]{article}
\usepackage[utf8]{inputenc}
\usepackage[italian]{babel}
\usepackage[margin=2.5cm]{geometry}
\usepackage{tikz}
\usepackage{booktabs}
\usepackage{array}
\usepackage{fancyhdr}
\usepackage{lastpage}
\usepackage{listings}
\usepackage{xcolor}

\lstset{
    basicstyle=\ttfamily\small,
    backgroundcolor=\color{gray!10},
    frame=single,
    breaklines=true,
    columns=fullflexible,
    keepspaces=true,
}

\usetikzlibrary{positioning, arrows.meta}

\setlength{\headheight}{14pt}
\pagestyle{fancy}
\fancyhf{}
\lhead{Esame Laboratorio Reti}
\rhead{A.A. 2025/2026}
\cfoot{Pagina \thepage\ di \pageref{LastPage}}

\begin{document}

\begin{center}
    {\LARGE \textbf{Esame di Laboratorio di Reti}}\\[0.3cm]
    {\large Anno Accademico 2025/2026}\\[0.3cm]
    \rule{\textwidth}{0.4pt}
\end{center}

\vspace{0.3cm}

\noindent
\begin{tabular}{@{}ll}
    \textbf{Nome e Cognome:} & \rule{8cm}{0.4pt} \\[0.3cm]
    \textbf{Matricola:}      & \rule{4cm}{0.4pt} \\[0.3cm]
    \textbf{Data:}           & 15 giugno 2026    \\
\end{tabular}

\vspace{0.5cm}

\noindent
\begin{minipage}[t]{0.5\textwidth}
    \textbf{Durata:} 1 ora e 30 minuti
\end{minipage}
\begin{minipage}[t]{0.5\textwidth}
    \raggedleft
    \begin{tabular}{|l|c|}
        \hline
        \textbf{Sezione}       & \textbf{Punti} \\
        \hline
        1. Configurazione rete & 2              \\
        2. Routing             & 6              \\
        3. Verifica            & 3              \\
        4. DNS                 & 2              \\
        \hline
        \textbf{Totale}        & \textbf{13}    \\
        \hline
    \end{tabular}
\end{minipage}

\vspace{0.5cm}

\section*{Scenario}

Configurare una rete composta da \textbf{5 router} (R1--R5) e \textbf{3 PC} (PC1--PC3), oltre a un server DNS. I quattro router esterni (R1--R4) sono collegati ad \textbf{anello} e ognuno è inoltre connesso al router centrale \textbf{R5}. Le rotte statiche devono utilizzare i \textbf{percorsi a costo minimo}.

\section*{Topologia}

\begin{center}
    \begin{tikzpicture}[
            router/.style={circle, draw, thick, minimum size=0.9cm, font=\small\bfseries},
            pc/.style={rectangle, draw, thick, minimum size=0.7cm, font=\small\bfseries},
            weight/.style={font=\scriptsize, fill=white, inner sep=1pt}
        ]

        % Routers: ring corners + central hub
        \node[router] (R1) at (0, 2.5) {R1};
        \node[router] (R2) at (6, 2.5) {R2};
        \node[router] (R3) at (6, -1.5) {R3};
        \node[router] (R4) at (0, -1.5) {R4};
        \node[router] (R5) at (3, 0.5) {R5};

        % Hosts
        \node[pc] (PC1) at (0, 4) {PC1};
        \node[pc] (PC2) at (6, 4) {PC2};
        \node[pc] (PC3) at (6, -3) {PC3};
        \node[pc] (DNS) at (0, -3) {DNS};

        % Host connections
        \draw[thick] (PC1) -- node[left, font=\scriptsize] {A} (R1);
        \draw[thick] (PC2) -- node[right, font=\scriptsize] {B} (R2);
        \draw[thick] (PC3) -- node[right, font=\scriptsize] {C} (R3);
        \draw[thick] (DNS) -- node[left, font=\scriptsize] {D} (R4);

        % Ring links
        \draw[thick] (R1) -- node[above, font=\scriptsize] {E} node[weight, below] {5} (R2);
        \draw[thick] (R2) -- node[right, font=\scriptsize] {F} node[weight, left] {1} (R3);
        \draw[thick] (R3) -- node[below, font=\scriptsize] {G} node[weight, above] {6} (R4);
        \draw[thick] (R4) -- node[left, font=\scriptsize] {H} node[weight, right] {2} (R1);

        % Spokes to hub R5
        \draw[thick] (R1) -- node[weight, pos=0.5] {J} node[weight, pos=0.7, above] {1} (R5);
        \draw[thick] (R2) -- node[weight, pos=0.5] {K} node[weight, pos=0.7, above] {3} (R5);
        \draw[thick] (R3) -- node[weight, pos=0.5] {L} node[weight, pos=0.7, below] {1} (R5);
        \draw[thick] (R4) -- node[weight, pos=0.5] {M} node[weight, pos=0.7, below] {2} (R5);

    \end{tikzpicture}
\end{center}

\section*{Schema indirizzi}

\begin{center}
    \begin{tabular}{cllc}
        \toprule
        \textbf{Dominio} & \textbf{Subnet} & \textbf{Collegamento}    & \textbf{Peso} \\
        \midrule
        A                & 10.0.1.0/24     & PC1 $\leftrightarrow$ R1 & --            \\
        B                & 10.0.2.0/24     & PC2 $\leftrightarrow$ R2 & --            \\
        C                & 10.0.3.0/24     & PC3 $\leftrightarrow$ R3 & --            \\
        D                & 10.0.4.0/24     & DNS $\leftrightarrow$ R4 & --            \\
        E                & 10.0.5.0/24     & R1 $\leftrightarrow$ R2  & 5             \\
        F                & 10.0.6.0/24     & R2 $\leftrightarrow$ R3  & 1             \\
        G                & 10.0.7.0/24     & R3 $\leftrightarrow$ R4  & 6             \\
        H                & 10.0.8.0/24     & R4 $\leftrightarrow$ R1  & 2             \\
        J                & 10.0.9.0/24     & R5 $\leftrightarrow$ R1  & 1             \\
        K                & 10.0.10.0/24    & R5 $\leftrightarrow$ R2  & 3             \\
        L                & 10.0.11.0/24    & R5 $\leftrightarrow$ R3  & 1             \\
        M                & 10.0.12.0/24    & R5 $\leftrightarrow$ R4  & 2             \\
        \bottomrule
    \end{tabular}
\end{center}

\section*{Consegne}

\subsection*{1. Configurazione di rete [2 punti]}

Scrivere il file \texttt{lab.conf} definendo tutte le macchine e i collision domain. Usare l'immagine \texttt{kathara/base}.

\subsection*{2. Configurazione routing [6 punti]}

Per ogni macchina, scrivere il file \texttt{.startup} con la configurazione degli indirizzi IP e le rotte statiche verso le sottoreti remote. I punti sono assegnati in base alla correttezza del funzionamento: tutti i PC e router devono poter comunicare tra loro.

\subsection*{3. Verifica percorsi ottimali [3 punti]}

Eseguire \texttt{traceroute} tra le coppie di PC e verificare che i percorsi siano ottimali (costo minimo).

\subsection*{4. Integrazione DNS [2 punti]}

Integrare il server DNS fornito (cartella \texttt{dns/}) collegandolo a R4 sul collision domain D con subnet 10.0.4.0/24.

Il server DNS è preconfigurato con gli indirizzi IP corretti dei PC (vedasi sotto).

Configurare i PC per utilizzare il DNS e verificare il funzionamento.

\subsection*{Convenzioni di indirizzamento (OBBLIGATORIE)}


\begin{itemize}
    \item \textbf{PC e DNS}: indirizzo che termina in \texttt{.1}
    \item \textbf{Router}: interfacce con indirizzo che termina in \texttt{.2} o \texttt{.3}
    \item \textbf{Sottoreti}: tutte le sottoreti devono essere implementate e raggiungibili, anche se non usate nella comunicazione tra pc
\end{itemize}

\textbf{Esempi:}
\begin{itemize}
    \item PC1 su subnet A: \texttt{10.0.1.1}, R1 su subnet A: \texttt{10.0.1.2}
    \item Su subnet E (R1$\leftrightarrow$R2): R1 usa \texttt{10.0.5.2}, R2 usa \texttt{10.0.5.3} (o viceversa)
\end{itemize}

\vfill

\section*{Riferimento comandi}

\textbf{Configurare un PC per usare il DNS} (da aggiungere nel .startup):
\begin{lstlisting}
echo "nameserver 10.0.4.1" > /etc/resolv.conf
\end{lstlisting}

\textbf{dns.startup}:
\begin{lstlisting}
ip address add 10.0.4.1/24 dev eth0
ip link set eth0 up
ip route add default via 10.0.4.2
systemctl start named
\end{lstlisting}

\textbf{dns/etc/bind/named.conf}:
\begin{lstlisting}
options {
    directory "/var/cache/bind";
    recursion yes;
    allow-query { any; };
    dnssec-validation no;
};

zone "lab.local" {
    type master;
    file "/etc/bind/db.lab.local";
};
\end{lstlisting}

\textbf{dns/etc/bind/db.lab.local}:
\begin{lstlisting}
$TTL    604800
@   IN  SOA  dns.lab.local. admin.lab.local. (
             1         ; Serial
             604800    ; Refresh
             86400     ; Retry
             2419200   ; Expire
             604800 )  ; Negative Cache TTL

@              IN  NS  dns.lab.local.
dns.lab.local. IN  A   10.0.4.1

pc1            IN  A   10.0.1.1
pc2            IN  A   10.0.2.1
pc3            IN  A   10.0.3.1
\end{lstlisting}

\end{document}
